% META: {"id":"1SPE-DERLOCAL-EX-010","chapitre":"1SPE-DERIVATION-LOCAL","type_objet":"exercice","capacites":["1SPE-DERIVATION-LOCAL-C2"],"capacites_codes":["C2"],"methodes":["M2"],"parcours":2,"competences":["calculer","raisonner"],"duree_min":12,"mode_creation":"ex_nihilo","sources_inspiration":[],"fichier_tex":"chapitres/1SPE-DERIVATION-LOCAL/exercices/1SPE-DERLOCAL-EX-010.tex","corrige_tex":"chapitres/1SPE-DERIVATION-LOCAL/corriges/1SPE-DERLOCAL-CO-010.tex","status":"approved","origin":"generated","status_history":["generated","audited","accepted","approved"],"release_acceptance":"RELEASE_OWNER_BATCH_ACCEPTANCE","acceptance_closure_digest":"sha256:9b3ccf9a81c5520fbb7e03b7b2d3e7bf2057a02834b6d908bf3be5f49f3b553f","acceptance_receipt_digest":"sha256:4fad86f0282e0fa4e0419cca1ad6379ae7af5a280c04a4a90880f90a1c044242"}
% BEGIN-VERIFY
% from sympy import symbols, simplify, Rational
% h = symbols('h', nonzero=True)
% # f(x) = x^3, f(1+h) = (1+h)^3 = 1+3h+3h^2+h^3
% taux = simplify(((1+h)**3 - 1)/h)
% assert simplify(taux - (h**2 + 3*h + 3)) == 0
% END-VERIFY
\begin{exercice}{1SPE-DERLOCAL-EX-010}{2}{12}
Soit $f(x)=x^3$.
\begin{enumerate}
  \item Développer $(1+h)^3$.
  \item En déduire l'expression de $\dfrac{f(1+h)-f(1)}{h}$ en fonction de $h$ pour $h
e 0$.
  \item Déterminer $f'(1)$.
  \item La tangente à la courbe de $f$ au point d'abscisse $1$ est-elle croissante ou décroissante ? Justifier.
\end{enumerate}
\end{exercice}
